\section{符号と道のり}\label{節:符号と道のり}
一直線上の運動でも,途中で向きが変わることがある.速度$v$は向きを含むので負にもなるが,\textbf{速さ}は$|v|$であり負にならない.

右向きに毎秒$2$メートルで$3$秒進み,次に左向きに毎秒$1$メートルで$2$秒進んだとする.位置の変化は
\[
 2\times3+(-1)\times2=4
\]
メートルである.一方,実際に通った道のりは
\[
 2\times3+1\times2=8
\]
メートルとなる.

一般の文字で書いておこう.正の時間$t_1$の間は速度$v_1$で,続く正の時間$t_2$の間は速度$v_2$で動いたとすれば,位置の変化と道のりはそれぞれ
\[
 v_1t_1+v_2t_2,\qquad |v_1|t_1+|v_2|t_2
\]
である.速度を時間について足せば位置の変化,速さを足せば道のりが得られるのだ.速度のグラフの下の面積として見るなら,位置の変化では$t$軸より下の部分を負の量として数え,道のりでは正の量として数えることになる.同じグラフの下の量でも,負の部分の数え方は求めるものによって変わるのである.
